\section{Gauge Fields, the Standard Model, and Generations}

Matter that only propagated and never interacted would not be our world. The same sites that carry the fermion also carry the forces between fermions, and they carry them without a second kind of stuff. This is where the substrate makes its boldest claim, that the gauge structure of the Standard Model is the symmetry of the $24$ directions, read at the cusp.

\subsection{One tone, and a photon}

The photon is not an addition. It is the vector octet, the $8v$ part of the sites, excited as a long-wavelength wave. Because that wave is charge-neutral it streams through the mesh unbound while charged knots stay localized, so the very same tone field that is matter when it is bound is light when it is waving. There is one substance and two motions, exactly as the monism promised. The wave carries a local $U(1)$ symmetry, the gauge freedom of electromagnetism, and coupling it to the charged fermion gives the Lorentz force and lattice quantum electrodynamics in the usual way. We tried at one point to make the photon a separate field, a second tone added alongside the first, and that route is recorded honestly among the abandoned ones in Appendix~\ref{app:abandoned}, because it violated the one-substance commitment and was not needed. The photon of the committed model is one tone all the way down.

That matter and light are the same field has a sharp, checkable consequence. The photon sector and the fermion sector are not two systems that happen to interact, they emerge together from one rule and feed back on each other both ways, the charge sourcing the field and the field deflecting the charge. The decisive test switches off the coupling and watches both effects die together, which they do, so the co-emergence is dynamical and measured, not a story told over two separate constructions (depth L2).

\subsection{Grand unification at the cusp}

The leap from electromagnetism to the full Standard Model is a leap in symmetry, and the symmetry is already present in $D_4$. Combined with the ternary tone, the root system furnishes the algebra $\mathrm{SO}(10)$, the grand-unified group that has long been the most economical home for one generation of matter. Its sixteen-dimensional spinor representation holds exactly the fifteen fermions of a single family plus a right-handed neutrino, with no room left over and none missing. The breaking of $\mathrm{SO}(10)$ down through $\mathrm{SU}(5)$ to the Standard Model is driven by a condensate of the self field, and the pattern of charges that results is the familiar one. These are group-theoretic facts about the representation, derived rather than simulated (depth: derived).

Grand unification makes quantitative predictions, and the substrate inherits them. The weak mixing angle comes out at $\sin^2\theta_W = 3/8$ at the unification scale, the standard $\mathrm{SU}(5)$ value, and running it down with the renormalization group lands near the measured $0.231$ for a spectrum with low-energy supersymmetry, which the model thereby predicts. The same structure gives bottom-tau mass unification and the tracelessness of the hypercharge, $\mathrm{Tr}\,Y = 0$, both at the unification scale, again as standard consequences. What is genuinely the substrate's contribution is not these relations, which belong to grand unification, but the claim that the group itself is forced by the geometry rather than chosen to fit.

\subsection{Three generations, and what is still open}

Matter comes in three families, and the substrate explains the count. Triality, the threefold symmetry among the octets, lifts to the exceptional Jordan algebra $J_3(\mathbb{O})$, the $3 \times 3$ Hermitian octonionic matrices, where the rank is forced to be three by the non-associativity of the octonions. Three is not fitted, it is the only rank the structure allows, and the three slots carry an exact $S_3$ family symmetry. This is the substrate's version of a result Boyle has pressed independently, and the count of three is solid (depth L1).

Here the honesty of the program shows. The three generation slots, as the symmetry leaves them, are \emph{degenerate}, they have the same mass. Splitting them into the three distinct masses we observe requires breaking the $S_3$ in a particular way, which is an open conjecture rather than a result, so the count is established and the splitting is not. The same line marks the strengths of the forces. The bare knit fixes the form of every dependence on a coupling, how a rate scales, where it vanishes, what it is proportional to, but it leaves the overall multiplicative value free. Deriving a number like the fine-structure constant $1/137$ needs a further principle, a fixed point or an anomaly condition or the full charged spectrum with its running, and none of that is in the bare rule. The mechanisms are in hand and the absolute numbers are not, a partial result reported as partial (depth L1), and gathered with the rest of the open frontier near the end.
